\documentclass[10pt,a4paper]{amsart}
\usepackage[margin=19mm]{geometry}
\usepackage{amsmath,amssymb,mathtools}
\usepackage[scr=rsfs,cal=euler]{mathalfa}
\usepackage{xfrac}
\usepackage[unicode,hidelinks,bookmarks=false]{hyperref}
\newcommand{\ie}{i.e.\ }
\newcommand{\cf}{cf.\ }
\newcommand{\resp}{respectively\ }
\newcommand{\Subsection}{\subsection}
\providecommand{\nobreakdash}{\nobreak\hbox{-}}
\setlength{\emergencystretch}{2em}
\multlinegap0pt
\usepackage{bm}
\usepackage{bbm}
\usepackage{dsfont}
\usepackage{xspace}
\usepackage{cancel}
\usepackage{mathtools}
\usepackage{cleveref}
\usepackage{amsthm}
\makeatletter
\@ifundefined{theorem}{\newtheorem{theorem}{Theorem}}{}
\@ifundefined{claim}{\newtheorem{claim}[theorem]{Claim}}{}
\@ifundefined{lemma}{\newtheorem{lemma}[theorem]{Lemma}}{}
\makeatother
\DeclareMathOperator{\surp}{surp}
\DeclareMathOperator{\trace}{tr}
\title[From small eigenvalues to large cuts, and Chowla's cosine pr]{From small eigenvalues to large cuts, and Chowla's cosine problem}
\author{Zhihan Jin, Aleksa Milojevi\'c, Istv\'an Tomon, Shengtong Zhang}
\date{Source: arXiv:2509.03490v2 (CC BY 4.0)}
\begin{document}
\maketitle

\noindent\emph{Source and licence.} From small eigenvalues to large cuts, and Chowla's cosine problem. Version arXiv:2509.03490v2 (CC BY 4.0), distributed under the Creative Commons Attribution 4.0 International licence, as recorded in the arXiv licence metadata for this exact version and shown on the abstract page https://arxiv.org/abs/2509.03490v2. Full rights notice: licenses/arxiv-2509.03490-CC-BY-4.0.txt.\par\smallskip

\noindent\textit{Editorial scope.} Counted unit: the Lemma whose source label is \texttt{lemma:main\_MaxCut} in the article (source0.tex lines 635--641, source label \texttt{lemma:main\_MaxCut}), delivered together with the article's own complete original proof of it (source0.tex lines 642--740, one \texttt{proof} environment of 99 lines). No mathematics is written, repaired or re-derived: statement and proof are the article's own text line for line. Required context retained uncounted: lemma:surp\_star (lines 415--422); lemma:smooth\_eigenvectors (lines 426--429); lemma:W-trace\_bound (lines 537--539); lemma:trace\_compression (lines 558--561); lemma:main\_smallest\_eigenalue (lines 571--576); claim:hadamard\_sum (lines 581--584). Every definition, display and label of those lines is retained unchanged. Omitted from this package and not redistributed: 1--414, 423--425, 430--536, 540--557, 562--570, 577--580, 585--634, 741--2486. Nothing inside a retained excerpt is omitted or altered: every retained body is delivered exactly as the article's lines are written, so no omission receipt of any kind accompanies this package. Citations: the retained excerpts carry no citation command at all, so no bibliography entry of the article is transcribed and no reference is redistributed. Reference adaptation: none was needed and none was made; every reference inside the delivered unit resolves inside it, so no unresolved reference or citation is produced. Printed slips kept literal: no printed word, symbol, hypothesis or number of the counted statement or of its proof was altered. Layout and class: the article's own class is \texttt{article} with packages including inputenc, t1enc, amsmath, amssymb, amsfonts, amsthm, amsopn, epsfig, hyphenat, float, graphicx, caption, appendix, epstopdf; the wrapper is the lane's pinned \texttt{amsart} editorial wrapper at 10pt on A4, which restates the theorem-like environments the retained text needs and adds the article's own macro definitions that the retained text needs (\texttt{\textbackslash surp}, \texttt{\textbackslash trace}), with extra wrapper packages (bm, bbm, dsfont, xspace, cancel, mathtools, cleveref). The delivered numbering is the unit's own. Assets: no retained span contains a figure environment, a graphics-inclusion command, a PDF-page inclusion, a table or a TikZ picture; no image file is copied into this package, the package declares no asset dependency and no image file is redistributed with it. Licence: version arXiv:2509.03490v2 of this article is distributed under the Creative Commons Attribution 4.0 International licence, as recorded in the arXiv licence metadata for this exact version and shown on the abstract page https://arxiv.org/abs/2509.03490v2. A full rights notice is delivered at licenses/arxiv-2509.03490-CC-BY-4.0.txt. Characters: every retained excerpt is pure ASCII, so the delivered engine, XeLaTeX with its default Latin Modern text font, reports no missing character.
\begin{lemma}\label{lemma:surp_star}
There exists an absolute constant $c>0$ such that the following holds. Let $G$ be a graph on $n$ vertices with eigenvalues $\lambda_i=\lambda_i(G)$, and let $\Delta^*=\Delta^*(G)$. Then
\begin{itemize}
    \item[(i)] $\surp^*(G)\geq \displaystyle\sum_{\lambda_i<0}|\lambda_i|$
    \item[(ii)] $\surp^*(G)\geq\frac{c}{\sqrt{\Delta^*+1}}\displaystyle\sum_{\lambda_i<0}\lambda_i^2$
    \item[(iii)]  $\surp^*(G)\geq \frac{c}{\Delta^*+1}\displaystyle\sum_{\lambda_i<0}|\lambda_i|^3$.
\end{itemize}
\end{lemma}

\begin{lemma}\label{lemma:smooth_eigenvectors}
Let $G$ be an $n$-vertex graph, and let $\lambda$ be an eigenvalue with normalized eigenvector $v$. Then
$$\|v\|_{\infty}\leq \frac{\sqrt{n}}{|\lambda|}.$$
\end{lemma}

\begin{lemma}\label{lemma:W-trace_bound}
$|\trace_W(M)|\leq \dim(W)^{1/2}\|M\|_F.$
\end{lemma}

\begin{lemma}\label{lemma:trace_compression}
Let $G$ be an $n$-vertex graph with adjacency matrix $A$ and let $W<\mathbb{R}^n$. Then for every $K>0$,
$$\trace_W(A)\leq S_K+K\dim(W).$$
\end{lemma}

\begin{lemma}\label{lemma:main_smallest_eigenalue}
Let $G$ be an $n$-vertex graph. If $T\geq 2|\lambda_n|\sqrt{n}$, then
\begin{equation}\label{eqn:S_T bound}
    4nS_{\frac{T^2}{2n}}\geq S_T^2.
\end{equation}
\end{lemma}

\begin{claim}\label{claim:hadamard_sum}
Let $v_1,\dots,v_n$ be an orthonormal basis of eigenvectors of $A$ corresponding to the eigenvalues $\lambda_1\geq \dots\geq \lambda_n$. Then:
    $$\sum_{\lambda_i,\lambda_j\geq T}\lambda_i\lambda_j\|v_i\circ v_j\|_2^2\geq\frac{S_T^2}{n}.$$
\end{claim}

\begin{lemma}\label{lemma:main_MaxCut}
For every $\gamma \in (0,1/6)$, there exists a constant $C>0$ such that the following holds.
Let $G$ be an $n$-vertex graph such that $\surp^*(G)\leq n^{1+\gamma}$. 
% Let $T\geq Cn^{1-\frac{1}{24}+\frac{\gamma}{4}}$ for some sufficiently large absolute constant $C>0$.
Then, for every $T\geq Cn^{1-\frac{1}{24}+\frac{\gamma}{4}}$,
$$4nS_{\frac{T^2}{2n}}\geq S_T^2.$$
\end{lemma}

\begin{proof}
Let $Q=\surp^*(G)$. 
As in the proof of \Cref{lemma:main_smallest_eigenalue}, we analyse the identity $A=A\circ A$. Let $E$ be the ``negative part'' of $A$, that is, $$E:=\sum_{\lambda_i<0}|\lambda_i|v_iv_i^T.$$ 
Then we can rewrite $A=A\circ A$ as
\begin{equation}\label{equ:mainlemma1}
A=(A+E)\circ (A+E)-2A\circ E-E\circ E.
\end{equation}
The proof revolves around choosing an appropriate subspace $W$ and bounding the $W$-traces of both sides. The terms $2A\circ E$ and $E\circ E$ constitute as error terms, for which we show that their contribution to the $W$-trace is negligible.

Let $W_0<\mathbb{R}^n$ be the subspace generated by those vectors $v_i\circ v_j$, where $\lambda_i$ and $\lambda_j$ are at least $T$, i.e.
$$W_0=\langle v_i\circ v_j : \lambda_i,\lambda_j\geq T\rangle.$$
The subspace $W_0$ is almost what we want. However, when bounding the error term $\trace_{W_0}(E\circ E)$, the large entries of $E$ have a non-negligible contribution. In order to overcome this, we introduce a cut-off 
$$\beta:=\frac{Q^{1/4}n^{7/8}}{T}>1.$$
Let $J\subset[n]$ be the set of indices $i$ such that $E_{i,i}>\beta$. Note that as $E$ is positive semidefinite, we have $\max_{i,j}|E_{i,j}|=\max_{i,i}E_{i,i}$, so $|E_{i,j}|\leq \beta$ for every $i,j\notin J$. Moreover, $|J|$ is small.



\begin{claim}
$|J|\leq Q/\beta.$
\end{claim}
\begin{proof}
    By Lemma \ref{lemma:surp_star} (i), we have 
        $\sum_{i=1}^n E_{i,i}=\trace(E)=\sum_{\lambda_i<0}|\lambda_i|\leq Q.$
    Hence, the sum $\sum_{i=1}^n E_{i,i}$ contains at most $Q/\beta$ terms larger than $\beta$.
\end{proof}

Let $Y<\mathbb{R}^n$ be the subspace of vectors that vanish on $J$, that is, 
$$Y:=\{y\in\mathbb{R}^n:\forall i\in J, y(i)=0\}.$$ Finally, define $$W:=\Pi_Y(W_0).$$
Note that $$\dim(W)\leq \dim(W_0)\leq N_T^2\leq \frac{S_T^2}{T^2}.$$
Consider the trace of the $W$-compressions of both sides of (\ref{equ:mainlemma1}). 
Let $K=\frac{T^2}{2n}$. 
Then \Cref{lemma:trace_compression} implies
$$\trace_W(A)\leq S_K+K\dim(W)\leq S_{K}+K\frac{S_T^2}{T^2}= S_{\frac{T^2}{2n}}+\frac{S_T^2}{2n}.$$
On the other hand, the term $\trace_W((A+E)\circ (A+E))$ can be lower bounded as follows.
\begin{claim}
    \[\trace_W((A+E)\circ (A+E))\geq \frac{S_T^2}{n}-\frac{Q  S_T^2 n^2}{\beta T^4}.\]    
\end{claim}
\begin{proof}
Since we have 
\begin{equation*}
(A+E)\circ (A+E)=  \bigg(\sum_{\lambda_i>0}\lambda_iv_iv_i^T\bigg)\circ \bigg(\sum_{\lambda_i>0}\lambda_iv_iv_i^T\bigg) =\sum_{\lambda_i,\lambda_j>0}\lambda_i\lambda_j (v_i\circ v_j)(v_i\circ v_j)^T,
\end{equation*}
we can write write
$$\trace_W((A+E)\circ (A+E))=\sum_{\lambda_i,\lambda_j>0} \lambda_i\lambda_j \|\Pi_W v_i\circ v_j\|_2^2\geq \sum_{\lambda_i,\lambda_j\geq T} \lambda_i\lambda_j \|\Pi_W v_i\circ v_j\|_2^2.$$
Here, by definition, we have $v_i\circ v_j\in W_0$, so $\Pi_W (v_i\circ v_j)=\Pi_Y (v_i\circ v_j)$. Thus,
$$\|\Pi_W (v_i\circ v_j)\|_2^2=\|\Pi_Y (v_i\circ v_j)\|_2^2=\|v_i\circ v_j\|_2^2-\sum_{k\in J} (v_i(k)v_j(k))^2.$$
By \Cref{lemma:smooth_eigenvectors}, the entries of $v_i$ and $v_j$ are bounded as $|v_i(k)|,|v_j(k)|\leq \frac{\sqrt{n}}{T}$, so we get
$$\big\|\Pi_W v_i\circ v_j \big\|_2^2
    =\|v_i\circ v_j\|_2^2-\sum_{k\in J} (v_i(k)v_j(k))^2
    \geq \|v_i\circ v_j\|_2^2-\frac{|J|n^2}{T^4}.$$
With this bound, we get 
\begin{align*}
    \trace_W((A+E)\circ (A+E))&\geq \sum_{\lambda_i,\lambda_j\geq T} \lambda_i\lambda_j\left(\|v_i\circ v_j\|_2^2-\frac{|J|n^2}{T^4}\right)\\
    &\geq \bigg(\,\sum_{\lambda_i,\lambda_j\geq T}\lambda_i\lambda_j\|v_i\circ v_j\|_2^2\,\bigg)-S_T^2\frac{|J|n^2}{T^4}
    \geq  \frac{S_T^2}{n}-\frac{Q  S_T^2 n^2}{\beta T^4}.
\end{align*}
Here, the second inequality follows by writing $\sum_{\lambda_i,\lambda_j\geq T}\lambda_i\lambda_j=S_T^2$, and the last inequality follows by Claim \ref{claim:hadamard_sum}, and writing  $|J|\leq Q/\beta$.
\end{proof}

Finally, we bound $\trace_W(E\circ A)$ and $\trace_W(E\circ E)$. First, we have 
$$\|E\circ A\|_F\leq \|E\|_F=\bigg(\sum_{\lambda_i<0}\lambda_i^2\bigg)^{1/2}=O(n^{1/4}Q^{1/2}),$$
where the last equality follows from \Cref{lemma:surp_star} (ii). 
By \Cref{lemma:W-trace_bound}, we have
$$\trace_W(A\circ E)=\dim(W)^{1/2} \|E\circ A\|_F\leq O\big(\dim(W)^{1/2}n^{1/4}Q^{1/2}\big).$$
Now consider $\trace_W(E\circ E)$.
Let $E'=E_Y$ be the $Y$-compression of $E$.
Then $E'_{i,j}=E_{i,j}$ if $i,j\not\in J$, and $E'_{i,j}=0$ otherwise. 
Recall that $|E_{i,j}| \le \beta$ for every $i,j \notin J$. We acquire
\begin{equation} \nonumber
    \|E'\circ E'\|_F
    =\bigg(\,\sum_{i,j\not\in J} E_{i,j}^4\,\bigg)^{1/2}
    \leq \beta \bigg(\,\sum_{i,j\not\in J}E_{i,j}^2\,\bigg)^{1/2}
    \leq \beta \|E\|_F=O(\beta n^{1/4}Q^{1/2}).
\end{equation}
From this, using \Cref{lemma:W-trace_bound} and that $W=\Pi_{Y} W_0$,
$$\trace_W(E\circ E)=\trace_W(E'\circ E')
    \leq \dim(W)^{1/2}\|E'\circ E'\|_F
    \leq O(\dim(W)^{1/2}\beta n^{1/4}Q^{1/2}).$$
Hence, the total contribution from the error terms can be bounded as
\begin{align*}
    2\trace_W(A\circ E)+\trace_W(E\circ E)
    &\leq O\big(\dim(W)^{1/2}\beta n^{1/4}Q^{1/2} \big)\\
    &\leq  O\big(\dim(W)\beta n^{1/4}Q^{1/2} \big)
    \leq O\!\left(\frac{S_T^2\beta n^{1/4}Q^{1/2}}{T^2}\right).
\end{align*}
In  the second inequality, we upper bounded $\dim(W)^{1/2}$ by $\dim(W)$, which is quite wasteful, but it simplifies upcoming calculations. Putting everything together, we proved that 
\begin{align*}
\trace_W(A\circ A)&=\trace_W((A+E)\circ (A+E))-2\trace_W(A\circ E)-\trace_W(E\circ E)\\
&\geq \frac{S_T^2}{n}-\frac{QS_T^2n^2}{\beta T^4}-O\left(\frac{S_T^2\beta n^{1/4}Q^{1/2}}{T^2}\right)=S_T^2\left(\frac{1}{n}-\frac{Qn^2}{\beta T^4}-O\left(\frac{\beta n^{1/4}Q^{1/2}}{T^2}\right)\right).
\end{align*}
The parameter $\beta$ was chosen such that the two negative terms have the same order of magnitude. After substituting $\beta=\frac{Q^{1/4}n^{7/8}}{T}$, we get
$$\trace_W(A\circ A) \ge S_T^2\left(\frac{1}{n}-O\left(\frac{Q^{3/4}n^{9/8}}{T^3}\right)\right)\geq \frac{3S_T^2}{4}.$$
Here, the last inequality holds by our assumptions that $Q\leq n^{1+\gamma}$ and $T\geq Cn^{1-\frac{1}{24}+\frac{\gamma}{4}}$. 
Now, comparing the left-hand-side and right-hand-side of (\ref{equ:mainlemma1}), we conclude the desired inequality by 
\begin{align}
S_{\frac{T^2}{2n}} + \frac{S_T^2}{2n} &\ge \trace_W(A) = \trace_W(A\circ A)
\geq \frac{3S_T^2}{4n}. \qedhere
\end{align}
\end{proof}

\end{document}
